% Options for packages loaded elsewhere
\PassOptionsToPackage{unicode}{hyperref}
\PassOptionsToPackage{hyphens}{url}
%
\documentclass[
  12pt,
]{article}
\usepackage{amsmath,amssymb}
\usepackage{iftex}
\ifPDFTeX
  \usepackage[T1]{fontenc}
  \usepackage[utf8]{inputenc}
  \usepackage{textcomp} % provide euro and other symbols
\else % if luatex or xetex
  \usepackage{unicode-math} % this also loads fontspec
  \defaultfontfeatures{Scale=MatchLowercase}
  \defaultfontfeatures[\rmfamily]{Ligatures=TeX,Scale=1}
\fi
\usepackage{lmodern}
\ifPDFTeX\else
  % xetex/luatex font selection
\fi
% Use upquote if available, for straight quotes in verbatim environments
\IfFileExists{upquote.sty}{\usepackage{upquote}}{}
\IfFileExists{microtype.sty}{% use microtype if available
  \usepackage[]{microtype}
  \UseMicrotypeSet[protrusion]{basicmath} % disable protrusion for tt fonts
}{}
\makeatletter
\@ifundefined{KOMAClassName}{% if non-KOMA class
  \IfFileExists{parskip.sty}{%
    \usepackage{parskip}
  }{% else
    \setlength{\parindent}{0pt}
    \setlength{\parskip}{6pt plus 2pt minus 1pt}}
}{% if KOMA class
  \KOMAoptions{parskip=half}}
\makeatother
\usepackage{xcolor}
\usepackage[margin=1in]{geometry}
\usepackage{graphicx}
\makeatletter
\def\maxwidth{\ifdim\Gin@nat@width>\linewidth\linewidth\else\Gin@nat@width\fi}
\def\maxheight{\ifdim\Gin@nat@height>\textheight\textheight\else\Gin@nat@height\fi}
\makeatother
% Scale images if necessary, so that they will not overflow the page
% margins by default, and it is still possible to overwrite the defaults
% using explicit options in \includegraphics[width, height, ...]{}
\setkeys{Gin}{width=\maxwidth,height=\maxheight,keepaspectratio}
% Set default figure placement to htbp
\makeatletter
\def\fps@figure{htbp}
\makeatother
\setlength{\emergencystretch}{3em} % prevent overfull lines
\providecommand{\tightlist}{%
  \setlength{\itemsep}{0pt}\setlength{\parskip}{0pt}}
\setcounter{secnumdepth}{5}
\usepackage{url}
\usepackage{setspace}
%\singlespacing
%\onehalfspacing
%\doublespacing
\usepackage{lineno}
%\linenumbers
\usepackage[belowskip=0pt,aboveskip=0pt]{caption}
\usepackage{relsize}
\newcommand{\Fmsy}{\ensuremath{F_{\text{MSY}}}\xspace}
\newcommand{\Fspr}[1]{\ensuremath{F_{\text{{#1}\%}}}\xspace}
\newcommand{\afrb}{Alaska Fishery Research Bulletin\xspace}
\newcommand{\ajms}{African Journal of Marine Science\xspace}
\newcommand{\amb}{Advances in Marine Biology\xspace}
\newcommand{\bms}{Bulletin of Marine Science\xspace}
\newcommand{\bjssf}{Bulletin of the Japanese Society of Scientific Fisheries\xspace}
\newcommand{\cb}{Conservation Biology\xspace}
\newcommand{\cjfas}{Canadian Journal of Fisheries and Aquatic Sciences\xspace}
\newcommand{\ea}{Ecological Applications\xspace}
\newcommand{\eer}{Evolutionary Ecology Research\xspace}
\newcommand{\elet}{Ecology Letters\xspace}
\newcommand{\emod}{Ecological Modelling\xspace}
\newcommand{\ebf}{Environmental Biology of Fishes\xspace}
\newcommand{\ff}{Fish and Fisheries\xspace}
\newcommand{\fo}{Fisheries Oceanography\xspace}
\newcommand{\fr}{Fisheries Research\xspace}
\newcommand{\fb}{Fishery Bulletin\xspace}
\newcommand{\ijms}{ICES Journal of Marine Science\xspace}
\newcommand{\iccat}{Collective Volume of Scientific Papers ICCAT\xspace}
\newcommand{\jae}{Journal of Animal Ecology\xspace}
\newcommand{\jai}{Journal of Applied Ichthyology\xspace}
\newcommand{\jdc}{Journal Du Conseil International Pour L'exploration De La Mer\xspace}
\newcommand{\jdcp}{Journal Du Conseil Permanent International Pour L'exploration De La Mer\xspace}
\newcommand{\jembe}{Journal of Experimental Marine Biology and Ecology\xspace}
\newcommand{\jfb}{Journal of Fish Biology\xspace}
\newcommand{\jsr}{Journal of Sea Research\xspace}
\newcommand{\jtb}{Journal of Theoretical Biology\xspace}
\newcommand{\jfrbc}{Journal of the Fisheries Research Board of Canada\xspace}
\newcommand{\jnwafs}{Journal of Northwest Atlantic Fisheries Science\xspace}
\newcommand{\mcf}{Marine and Coastal Fisheries: Dynamics, Management, and Ecosystem Science\xspace}
\newcommand{\mb}{Marine Biology\xspace}
\newcommand{\meps}{Marine Ecology Progress Series\xspace}
\newcommand{\mfr}{Marine Fisheries Review\xspace}
\newcommand{\mpb}{Marine Pollution Bulletin\xspace}
\newcommand{\najfm}{North American Journal of Fisheries Management\xspace}
\newcommand{\nzjmfr}{New Zealand Journal of Marine and Freshwater Research\xspace}
\newcommand{\pnas}{Proceedings of the National Academy of Sciences USA\xspace}
\newcommand{\rpvrciemm}{Rapports et Proc\`es-Verbaux des R\'eunions. Conseil Internationale pour l'Exploration de la Mer\xspace}
\newcommand{\rpvrcpiemm}{Rapports et Proc\`es-Verbaux des R\'eunions. Conseil Permanent Internationale pour l'Exploration de la Mer\xspace}
\newcommand{\rfbf}{Reviews in Fish Biology and Fisheries\xspace}
\newcommand{\sajms}{South African Journal of Marine Science\xspace}
\newcommand{\tafs}{Transactions of the American Fisheries Society\xspace}

\newcommand{\anzjs}{Australian \& New Zealand Journal of Statistics\xspace}
\newcommand{\as}{Applied Statistics\xspace}
\newcommand{\csda}{Computational Statistics \& Data Analysis\xspace}
\newcommand{\ees}{Environmental and Ecological Statistics\xspace}
\newcommand{\jas}{Journal of Applied Statistics\xspace}
\newcommand{\jabes}{Journal of Agricultural, Biological, and Environmental Statistics\xspace}
\newcommand{\jasa}{Journal of the American Statistical Association\xspace}
\newcommand{\jrssb}{Journal of the Royal Statistical Society. Series B\xspace}
\newcommand{\sm}{Statistics in Medicine}

\usepackage{float}
\usepackage{amsmath}
\usepackage{pdflscape}
\usepackage{hyperref}
\usepackage{xspace}
\usepackage{bm}
\usepackage{caption,graphics}
\usepackage{graphicx}
\usepackage{makecell}
\usepackage{lineno}
\linenumbers
\renewcommand\figurename{Fig.}
\renewcommand{\figureautorefname}{Fig.}
\captionsetup{labelsep=period, singlelinecheck=false}
\newcommand{\changesize}[1]{\fontsize{#1pt}{#1pt}\selectfont}
\renewcommand{\arraystretch}{1.5}
\renewcommand\theadfont{}
\usepackage{booktabs}
\usepackage{longtable}
\usepackage{array}
\usepackage{multirow}
\usepackage{wrapfig}
\usepackage{float}
\usepackage{colortbl}
\usepackage{pdflscape}
\usepackage{tabu}
\usepackage{threeparttable}
\usepackage{threeparttablex}
\usepackage[normalem]{ulem}
\usepackage{makecell}
\usepackage{xcolor}
\ifLuaTeX
  \usepackage{selnolig}  % disable illegal ligatures
\fi
\usepackage[]{natbib}
\bibliographystyle{cjfas2.bst}
\IfFileExists{bookmark.sty}{\usepackage{bookmark}}{\usepackage{hyperref}}
\IfFileExists{xurl.sty}{\usepackage{xurl}}{} % add URL line breaks if available
\urlstyle{same}
\hypersetup{
  pdftitle={Random effects on numbers-at-age transitions implicitly account for movement dynamics and improve stock assessment and management},
  pdfauthor={Chengxue Li\^{}\{1,6,7,*\}; Jonathan J. Deroba1; Aaron M. Berger2; Daniel R. Goethel3; Brian J. Langseth4; Amy M. Schueller5; Timothy J. Miller1},
  hidelinks,
  pdfcreator={LaTeX via pandoc}}

\title{Random effects on numbers-at-age transitions implicitly account
for movement dynamics and improve stock assessment and management}
\author{Chengxue Li\(^{1,6,7,*}\) \and Jonathan J.
Deroba\textsuperscript{1} \and Aaron M.
Berger\textsuperscript{2} \and Daniel R.
Goethel\textsuperscript{3} \and Brian J.
Langseth\textsuperscript{4} \and Amy M.
Schueller\textsuperscript{5} \and Timothy J. Miller\textsuperscript{1}}
\date{}

\begin{document}
\maketitle

\(^{1}\) NOAA, Northeast Fisheries Science Center, 166 Water Street,
Woods Hole, MA 02543, USA

\(^{2}\) NOAA, Northwest Fisheries Science Center, 2032 SE OSU Drive,
Newport, OR 97365, USA

\(^{3}\) NOAA, Alaska Fisheries Science Center, 17109 Point Lena Loop
Road, Juneau, AK 99801, USA

\(^{4}\) NOAA, Northwest Fisheries Science Center, 2725 Montlake Blvd E,
Seattle, WA 98112, USA

\(^{5}\) NOAA, Southeast Fisheries Science Center, 101 Pivers Island
Road, Beaufort, NC 28516, USA

\(^{6}\) School of Marine and Atmospheric Sciences, Stony Brook
University, Stony Brook, NY 11794, USA

\(^{7}\) Saltwater Inc., 733 N Street, Anchorage, AK 99501, USA

\(^{*}\) Corresponding author:
\href{mailto:chengxue.li@noaa.gov}{\nolinkurl{chengxue.li@noaa.gov}}

\vspace{0.5cm}

\textbf{ORCID:}

Chengxue Li:
\href{https://orcid.org/0000-0001-7518-7234}{0000-0001-7518-7234}

Aaron M. Berger:
\href{https://orcid.org/0000-0002-1408-7122}{0000-0002-1408-7122}

Daniel R. Goethel
\href{https://orcid.org/0000-0001-7121-6181}{0000-0001-7121-6181}

Brian J. Langseth
\href{https://orcid.org/0000-0002-9901-6146}{0000-0002-9901-6146}

Timothy J. Miller:
\href{https://orcid.org/0000-0003-1411-1206}{0000-0003-1411-1206}

\pagebreak

\hypertarget{methods}{%
\subsection*{Methods}\label{methods}}
\addcontentsline{toc}{subsection}{Methods}

\hypertarget{population-dynamics}{%
\subsubsection*{Population Dynamics}\label{population-dynamics}}
\addcontentsline{toc}{subsubsection}{Population Dynamics}

For population \(s\), the number-at-age \(a\) that survived
(\(N_{S,s,a,t+\delta}\)), were harvested from the fishery
(\(N_{H,s,a,t+\delta}\)), and died from natural mortality
(\(N_{D,s,a,t+\delta}\)) over the time interval \(\delta\) are given by:

\setcounter{equation}{0}

\begin{equation}
\begin{aligned}
N_{S,s,a,t+\delta} &= N_{s,a,t} S(s,a,\delta) = N_{s,a,t} e^{-Z_{s,a}\delta} \\
N_{H,s,a,t+\delta} &= N_{s,a,t} H(s,a,\delta) = N_{s,a,t} \frac{F_{s,a}}{Z_{s,a}}\left(1 - e^{-Z_{s,a}\delta}\right) \\
N_{D,s,a,t+\delta} &= N_{s,a,t} D(s,a,\delta) = N_{s,a,t} \frac{M_{s,a}}{Z_{s,a}}\left(1 - e^{-Z_{s,a}\delta}\right)
\end{aligned}
\end{equation}

where \(N_{s,a,t}\) is the number-at-age \(a\) at time \(t\) for
population \(s\); \(S\), \(H\), and \(D\) represent the proportions
(fractions) of individuals that survived, were harvested, or died from
natural mortality, respectively, over the time interval \(\delta\);
\(F_{t,a}\) and \(M_{t,a}\) denote instantaneous rates of fishing
mortality and natural mortality at age \(a\) and time \(t\), and
\(Z_{s,a}\) denotes total mortality rate, where
\(Z_{s,a} = F_{s,a} + M_{s,a}\).

The probability transition matrix, when there is movement between
regions, can be generalized as:

\begin{equation}
\mathbf{P}_{s,a,t,\delta} = \begin{bmatrix}
 \mathbf{O}(s,a,\delta) & \mathbf{H}(s,a,\delta) & \mathbf{D}(s,a,\delta) \\
 0 & \mathbf{I}_{H} & 0 \\
 0 & 0 & \mathbf{I}_{D}
\end{bmatrix}
\end{equation}

where \(\mathbf{O}\) is the \(n_R \times n_R\) matrix representing the
combined probability of survival and movement between regions. Assuming
movement occurs sequentially after survival, \(\mathbf{O}\) is the
product of the survival probability matrix (\(\mathbf{S}\)) and movement
probability matrix \(\mu\), such that \(\mathbf{O} = \mathbf{S} \mu\).
\(\mathbf{H}\) is the \(n_R \times n_F\) matrix representing proportions
of individuals harvested by each fleet in each region, \(\mathbf{D}\) is
the \(n_R \times n_R\) matrix representing proportions of individuals
that died from natural mortality within each region, and
\(\mathbf{I}_{H}\) and \(\mathbf{I}_{D}\), which have the same
dimensions as \(\mathbf{H}\) and \(\mathbf{D}\)), are identity matrices
that ensure the overall dimensions of the probability transition matrix
remain consistent when combined. If there is only one region and one
fleet, these identity matrices reduce to scalars of 1.

For population \(s\), the total number-at-age \(a\) at the end of time
interval \(\delta\) can be generalized as:

\begin{equation}
\mathbf{N}_{s,a,t+\delta} = \mathbf{P}(s,a,\delta)' \mathbf{N}_{s,a,t}
\end{equation}

where the numbers of individuals that survived, were harvested, and died
from natural mortality are decomposed as:

\begin{equation}
\begin{aligned}
\mathbf{N}_{S,s,a,t+\delta} &= \mathbf{O}(s,a,\delta)' \mathbf{N}_{s,a,t} \\
\mathbf{N}_{H,s,a,t+\delta} &= \mathbf{H}(s,a,\delta)' \mathbf{N}_{s,a,t} \\
\mathbf{N}_{D,s,a,t+\delta} &= \mathbf{D}(s,a,\delta)' \mathbf{N}_{s,a,t}
\end{aligned}
\end{equation}

\hypertarget{population-structure-and-seasonality}{%
\subsubsection*{Population Structure and
Seasonality}\label{population-structure-and-seasonality}}
\addcontentsline{toc}{subsubsection}{Population Structure and
Seasonality}

A natal homing population structure is assumed in WHAM. Each population
is associated with a natal region, and all fish originating in that
region must return there for spawning. Outside of the spawning season,
fish can move between regions (see details in the sections on OM or EM
movement dynamics). Seasonality is included to support the assumption of
natal homing. Specifically, a number of time intervals within a year
(\(\Sigma_{i=1}^{n} \delta_i = 1\)) are used to indicate movement
dynamics prior to the spawning season, during the spawning season, and
after the spawning season. The probability transition matrix for the
entire year is a product of all ``seasonal'' matrices:

\begin{equation}
\mathbf{P}_{s,a,y} = \prod^{n}_{t=1} \mathbf{P}_{s,a,y,t}
\end{equation}

\hypertarget{fishing-mortality-f}{%
\subsubsection*{\texorpdfstring{Fishing Mortality
(\(F\))}{Fishing Mortality (F)}}\label{fishing-mortality-f}}
\addcontentsline{toc}{subsubsection}{Fishing Mortality (\(F\))}

\({F}\) at age \({a}\) for fleet \({f}\) is calculated as a product of
fully selected \({F}\) and selectivity-at-age (i.e., the separability
assumption) for fleet \({f}\):
\(\tilde F_{f,a} = \tilde F \textit{Sel}_{f,a}\). In WHAM, a fully
selected, total \({F}\) is the maximum of the total \({F}\)-at-age
summed across all fleets:

\begin{equation}
\begin{aligned}
F_{\text{total},a,y} &= \sum_{i=1}^{n_f} F_{f_i,a,y} \\
F_{\text{total},y} &= \operatorname*{argmax}_a F_{\text{total},a,y} \\
\end{aligned}
\end{equation}

\hypertarget{reference-points}{%
\subsubsection*{Reference Points}\label{reference-points}}
\addcontentsline{toc}{subsubsection}{Reference Points}

The fully selected total \(F\) defined above was used to estimate a
single universal \({F}\) reference point applicable to all populations
and regions. A `static' SPR-based reference point was calculated by
averaging all of the inputs (e.g., selectivity, natural mortality,
movement, weight-at-age, maturity-at-age) over the most recent 5 model
years for each region, except for recruitment, which was averaged over
all model years for each region.

For \(X\%\) SPR-based reference points, WHAM uses a Newton method to
solve for the universal \(F_{X\%}\) reference point:

\begin{equation}
\tilde{F}^{(i)} = \tilde{F}^{(i-1)} - \frac{g\left(\tilde{F}^{(i-1)}\right)}{\tilde{F}^{(i-1)}g'\left(\tilde{F}^{(i-1)}\right)}
\end{equation}

where \(g(F)\) is the difference between the weighted spawning biomass
per recruit at \(F\) and \(X\%\) of the weighted unfished spawning
biomass per recruit:

\begin{equation}
g(F) = \sum_{s=1}^S w_s \left( \frac{\tilde{\tilde{SSB}}}{\tilde{R}(F)} - \frac{X}{100} \frac{\tilde{\tilde{SSB}}}{\tilde{R}(F=0)} \right)
\end{equation}

Here weights are determined by the average recruitment for different
regions/populations:

\begin{equation}
w_i = \frac{\bar{R}_i}{\sum_{i=1}^{n_s} \bar{R}_i}
\end{equation}

The fully selected total \(F\) reference point can be disaggregated to
regional \(F\) reference points, given the relative contribution of
selectivity (\(Sel\)) at age \(a\) from each fleet in each region \(R\):

\begin{equation}
F_{X\%_{\text{a},R}} = F_{X\%} \text{$Sel$}_{a,R}
\end{equation}

The regional \(F_{X\%}\) (e.g., \(F_{40\%}\)) does not guarantee
achieving 40\% of region-specific SPR, but serve to collectively achieve
40\% of SPR at the global scale. Regional \(F\) reference points are not
currently implemented, as a unique solution is not guaranteed in the
presence of movement between regions. An analogous Newton method is used
to solve for \(F\) that maximizes yield for MSY-based reference points.

\pagebreak

\renewcommand{\thetable}{S\arabic{table}}
\setcounter{table}{0}

\begin{table}[H]
    \centering
    \caption{Biological reference point (BRP) and harvest control rule (HCR) used for 10 estimation models.}
    \label{HCR_table}
    \begin{tabular}{p{2cm} p{2cm} p{2cm} p{8cm}}
        \toprule
        \textbf{Model} & \textbf{Global BRP} & \textbf{Rec Weights} & \textbf{Catch Apportionment} \\
        \midrule
        PAN$_{\text{NAA}}$ & Yes & No & Regional survey biomass \\
        PAN$_{\text{noNAA}}$ & Yes & No & Regional survey biomass \\
        FAA$_{\text{NAA}}$ & Yes & No & Disaggregated regional $F$ from global $F$\\
        FAA$_{\text{noNAA}}$ & Yes & No & Disaggregated regional $F$ from global $F$ \\
        SEP$_{\text{NAA}}$ & No & No & Independent regional $F$ \\
        SEP$_{\text{noNAA}}$ & No & No & Independent regional $F$ \\
        SpD$_{\text{NAA}}$ & Yes & Yes & Disaggregated regional $F$ from global $F$ \\
        SpD$_{\text{noNAA}}$ & Yes & Yes & Disaggregated regional $F$ from global $F$ \\
        SpE$_{\text{NAA,Est}}$ & Yes & Yes & Disaggregated regional $F$ from global $F$ \\
        SpE$_{\text{NAA,Fix}}$ & Yes & Yes & Disaggregated regional $F$ from global $F$ \\
        \bottomrule
    \end{tabular}
\end{table}

\renewcommand{\thetable}{S\arabic{table}}

\renewcommand{\thefigure}{S\arabic{figure}}
\setcounter{figure}{0}

\begin{figure}[H]
\centering
\includegraphics[width=\textwidth]{Movement1.PNG}
\caption{Simulated AR1 movement rates (from one replicate) under the high movement scenario ($\text{OM}_{\text{high}}$). The panels show time‐varying movement rates between regions, with green indicating movement from region 1 to region 2 and red indicating movement from region 2 to region 1. The same movement rates were applied in both the spring and winter seasons.}
\label{fig:supp_movement1}
\end{figure}

\begin{figure}[H]
\centering
\includegraphics[width=\textwidth]{Bias_Rec_OM_1.PNG}
\caption{Relative bias in recruitment for region 1, region 2, and global recruitment from the first assessment model during the feedback period under the high movement scenario ($\text{OM}_{\text{high}}$). Each point represents the median across the 100 realizations, and error bars represent the interquartile range (25th–75th percentiles). EMs with a grey background indicate that NAA random effects were included, while EMs with a white background indicate that NAA random effects were not included. Three fishing scenarios are included: Adaptive ($\mathrm{F_A}$), Unsustainable ($\mathrm{F_U}$), and Contrasting ($\mathrm{F_C}$).}
\label{fig:supp_Bias_Rec_OM_1}
\end{figure}

\begin{figure}[H]
\centering
\includegraphics[width=\textwidth]{Bias_SSB_OM_1.PNG}
\caption{Relative bias in $SSB$ for region 1, region 2, and global $SSB$ from the first assessment model during the feedback period under the high movement scenario ($\text{OM}_{\text{high}}$). Each point represents the median across the 100 realizations, and error bars represent the interquartile range (25th–75th percentiles). EMs with a grey background indicate that NAA random effects were included, while EMs with a white background indicate that NAA random effects were not included. Three fishing scenarios are included: Adaptive ($\mathrm{F_A}$), Unsustainable ($\mathrm{F_U}$), and Contrasting ($\mathrm{F_C}$).}
\label{fig:supp_Bias_SSB_OM_1}
\end{figure}

\begin{figure}[H]
\centering
\includegraphics[width=\textwidth]{Parameter_Estimates_OM_1.PNG}
\caption{The mean recruitment parameter, recruitment standard deviation, and NAA standard deviation from the last assessment model during the feedback period under the high movement scenario ($\text{OM}_{\text{high}}$). Each point represents the median across the 100 realizations, and error bars represent the interquartile range (25th–75th percentiles). EMs with a grey background indicate that NAA random effects were included, while EMs with a white background indicate that NAA random effects were not included. Three fishing scenarios are included: Adaptive ($\mathrm{F_A}$), Unsustainable ($\mathrm{F_U}$), and Contrasting ($\mathrm{F_C}$).}
\label{fig:supp_Parameter_Estimates_OM_1}
\end{figure}

\begin{figure}[H]
\centering
\includegraphics[width=\textwidth]{Catch_G_OM1.PNG}
\caption{Relative difference in annual total catch between the baseline EM ($\text{SpE}_{\scriptscriptstyle \text{NAA,Fix}}$) and other EMs for each fishing history under the high movement scenario ($\text{OM}_{\text{high}}$). Each line represents the median trajectory across the 100 realizations, and the shaded area represents the interquartile range (40th–60th percentiles). The performance of EMs with and without NAA random effects is shown separately for better comparison. Three fishing scenarios are included: Adaptive ($\mathrm{F_A}$), Unsustainable ($\mathrm{F_U}$), and Contrasting ($\mathrm{F_C}$).}
\label{fig:supp_annual_catch_G_OM1}
\end{figure}

\begin{figure}[H]
\centering
\includegraphics[width=\textwidth]{Catch_1_OM1.PNG}
\caption{Relative difference in annual catch in region 1 between the baseline EM ($\text{SpE}_{\scriptscriptstyle \text{NAA,Fix}}$) and other EMs for each fishing history under the high movement scenario ($\text{OM}_{\text{high}}$). Each line represents the median trajectory across the 100 realizations, and the shaded area represents the interquartile range (40th–60th percentiles). The performance of EMs with and without NAA random effects is shown separately for better comparison. Three fishing scenarios are included: Adaptive ($\mathrm{F_A}$), Unsustainable ($\mathrm{F_U}$), and Contrasting ($\mathrm{F_C}$).}
\label{fig:supp_annual_catch_1_OM1}
\end{figure}

\begin{figure}[H]
\centering
\includegraphics[width=\textwidth]{Catch_2_OM1.PNG}
\caption{Relative difference in annual catch in region 2 between the baseline EM ($\text{SpE}_{\scriptscriptstyle \text{NAA,Fix}}$) and other EMs for each fishing history under the high movement scenario ($\text{OM}_{\text{high}}$). Each line represents the median trajectory across the 100 realizations, and the shaded area represents the interquartile range (40th–60th percentiles). The performance of EMs with and without NAA random effects is shown separately for better comparison. Three fishing scenarios are included: Adaptive ($\mathrm{F_A}$), Unsustainable ($\mathrm{F_U}$), and Contrasting ($\mathrm{F_C}$).}
\label{fig:supp_annual_catch_2_OM1}
\end{figure}

\begin{figure}[H]
\centering
\includegraphics[width=\textwidth]{SSB_G_OM1.PNG}
\caption{Relative difference in annual total $SSB$ between the baseline EM ($\text{SpE}_{\scriptscriptstyle \text{NAA,Fix}}$) and other EMs for each fishing history under the high movement scenario ($\text{OM}_{\text{high}}$). Each line represents the median trajectory across the 100 realizations, and the shaded area represents the interquartile range (40th–60th percentiles). The performance of EMs with and without NAA random effects is shown separately for better comparison. Three fishing scenarios are included: Adaptive ($\mathrm{F_A}$), Unsustainable ($\mathrm{F_U}$), and Contrasting ($\mathrm{F_C}$).}
\label{fig:supp_SSB_G_OM1}
\end{figure}

\begin{figure}[H]
\centering
\includegraphics[width=\textwidth]{SSB_1_OM1.PNG}
\caption{Relative difference in annual $SSB$ for region 1 between the baseline EM ($\text{SpE}_{\scriptscriptstyle \text{NAA,Fix}}$) and other EMs for each fishing history under the high movement scenario ($\text{OM}_{\text{high}}$). Each line represents the median trajectory across the 100 realizations, and the shaded area represents the interquartile range (40th–60th percentiles). The performance of EMs with and without NAA random effects is shown separately for better comparison. Three fishing scenarios are included: Adaptive ($\mathrm{F_A}$), Unsustainable ($\mathrm{F_U}$), and Contrasting ($\mathrm{F_C}$).}
\label{fig:supp_SSB_1_OM1}
\end{figure}

\begin{figure}[H]
\centering
\includegraphics[width=\textwidth]{SSB_2_OM1.PNG}
\caption{Relative difference in annual $SSB$ for region 2 between the baseline EM ($\text{SpE}_{\scriptscriptstyle \text{NAA,Fix}}$) and other EMs for each fishing history under the high movement scenario ($\text{OM}_{\text{high}}$). Each line represents the median trajectory across the 100 realizations, and the shaded area represents the interquartile range (40th–60th percentiles). The performance of EMs with and without NAA random effects is shown separately for better comparison. Three fishing scenarios are included: Adaptive ($\mathrm{F_A}$), Unsustainable ($\mathrm{F_U}$), and Contrasting ($\mathrm{F_C}$).}
\label{fig:supp_SSB_2_OM1}
\end{figure}

\begin{figure}[H]
\centering
\includegraphics[width=\textwidth]{Probability_Overfished_Overfishing_OM1_SPR.PNG}
\caption{Probability of $SSB < SSB_{40\%}$ and $F > F_{40\%}$ for each EM over the 30-year feedback period under the high movement scenario ($\text{OM}_{\text{high}}$). Each point represents the median across the 100 realizations (one value per realization), and error bars represent the interquartile range (25th–75th percentiles). EMs with a grey background indicate that NAA random effects were included, while EMs with a white background indicate that NAA random effects were not included. Three fishing scenarios are included: Adaptive ($\mathrm{F_A}$), Unsustainable ($\mathrm{F_U}$), and Contrasting ($\mathrm{F_C}$).}
\label{fig:Overfished_OM1_SPR}
\end{figure}

\begin{figure}[H]
\centering
\includegraphics[width=\textwidth]{Population_Status_OM_1_MSY.PNG}
\caption{Terminal-year status of $F$ and $SSB$ at the global scale for each EM under the high movement scenario ($\text{OM}_{\text{high}}$). Each point represents the median across the 100 realizations, and error bars represent the interquartile range (25th–75th percentiles). The vertical dashed line indicates the threshold of overfished ($SSB_{T}/SSB_{MSY_{T}} < 0.5$), and the horizontal dashed line indicates the threshold of overfishing ($F_{T}/F_{MSY_{T}} > 1$). The performance of EMs with and without NAA random effects is shown separately for better comparison. Three fishing scenarios are included: Adaptive ($\mathrm{F_A}$), Unsustainable ($\mathrm{F_U}$), and Contrasting ($\mathrm{F_C}$).}
\label{fig:supp_population_status_G_OM_1_MSY}
\end{figure}

\begin{figure}[H]
\centering
\includegraphics[width=\textwidth]{Population_Status_OM_1_SPR.PNG}
\caption{Terminal-year status of $F$ and $SSB$ at the global scale for each EM under the high movement scenario ($\text{OM}_{\text{high}}$). Each point represents the median across the 100 realizations, and error bars represent the interquartile range (25th–75th percentiles). The vertical dashed line indicates the threshold of overfished ($SSB_{T}/SSB_{40\%_{T}} < 0.5$), and the horizontal dashed line indicates the threshold of overfishing ($F_{T}/F_{40\%_{T}} > 1$). The performance of EMs with and without NAA random effects is shown separately for better comparison. Three fishing scenarios are included: Adaptive ($\mathrm{F_A}$), Unsustainable ($\mathrm{F_U}$), and Contrasting ($\mathrm{F_C}$).}
\label{fig:supp_population_status_G_OM_1_SPR}
\end{figure}

\begin{figure}[H]
\centering
\includegraphics[width=\textwidth,height=\textheight]{Holistic_OM_R1_1.PNG}
\caption{Overall performance of each EM in region 1 under the high movement scenario ($\text{OM}_{\text{high}}$), including the following relative performance metrics: 1) long-term catch ($\text{Catch}_{\text{long}}$); 2) short-term catch ($\text{Catch}_{\text{short}}$); 3) long-term $SSB$ ($SSB_{\text{long}}$); 4) short-term $SSB$ ($SSB_{\text{short}}$); 5) probability of $F < F_{MSY}$ ($\text{P}_{F<F_{MSY}}$); 6) probability of $SSB > SSB_{MSY}$ ($\text{P}_{SSB>SSB_{MSY}}$); 7) overfishing status in the terminal year ($\text{OFG}_{\text{terminal}}$); 8) overfished status in the terminal year ($\text{OFD}_{\text{terminal}}$); and 9) average annual catch variation (AACV). Bars represent the median score across the 100 realizations, and error bars represent the interquartile range (25th–75th percentiles). The number displayed in each panel represents the total score (sum of all medians for that EM) for the corresponding EM under that fishing scenario, with higher totals indicating better overall performance. The black dashed line separates the performance of each EM. EMs with a grey background indicate that NAA random effects were included, while EMs with a white background indicate that NAA random effects were not included. Three fishing scenarios are included: Adaptive ($\mathrm{F_A}$), Unsustainable ($\mathrm{F_U}$), and Contrasting ($\mathrm{F_C}$).}
\label{fig:supp_holistic1_R1}
\end{figure}

\begin{figure}[H]
\centering
\includegraphics[width=\textwidth,height=\textheight]{Holistic_OM_R2_1.PNG}
\caption{Overall performance of each EM in region 2 under the high movement scenario ($\text{OM}_{\text{high}}$), including the following relative performance metrics: 1) long-term catch ($\text{Catch}_{\text{long}}$); 2) short-term catch ($\text{Catch}_{\text{short}}$); 3) long-term $SSB$ ($SSB_{\text{long}}$); 4) short-term $SSB$ ($SSB_{\text{short}}$); 5) probability of $F < F_{MSY}$ ($\text{P}_{F<F_{MSY}}$); 6) probability of $SSB > SSB_{MSY}$ ($\text{P}_{SSB>SSB_{MSY}}$); 7) overfishing status in the terminal year ($\text{OFG}_{\text{terminal}}$); 8) overfished status in the terminal year ($\text{OFD}_{\text{terminal}}$); and 9) average annual catch variation (AACV). Bars represent the median score across the 100 realizations, and error bars represent the interquartile range (25th–75th percentiles). The number displayed in each panel represents the total score (sum of all medians for that EM) for the corresponding EM under that fishing scenario, with higher totals indicating better overall performance. The black dashed line separates the performance of each EM. EMs with a grey background indicate that NAA random effects were included, while EMs with a white background indicate that NAA random effects were not included. Three fishing scenarios are included: Adaptive ($\mathrm{F_A}$), Unsustainable ($\mathrm{F_U}$), and Contrasting ($\mathrm{F_C}$).}
\label{fig:supp_holistic1_R2}
\end{figure}

\% --- Low movement section continues here ---

\begin{figure}[H]
\centering
\includegraphics[width=\textwidth]{Movement2.PNG}
\caption{Simulated AR1 movement rates (from one replicate) under the low movement scenario ($\text{OM}_{\text{low}}$). The panels show time‐varying movement rates between regions, with green indicating movement from region 1 to region 2 and red indicating movement from region 2 to region 1. The same movement rates were applied in both the spring and winter seasons.}
\label{fig:supp_movement2}
\end{figure}

\begin{figure}[H]
\centering
\includegraphics[width=\textwidth]{Bias_Rec_OM_2.PNG}
\caption{Relative bias in recruitment for region 1, region 2, and global recruitment from the first assessment model during the feedback period under the low movement scenario ($\text{OM}_{\text{low}}$). Each point represents the median across the 100 realizations, and error bars represent the interquartile range (25th–75th percentiles). EMs with a grey background indicate that NAA random effects were included, while EMs with a white background indicate that NAA random effects were not included. Three fishing scenarios are included: Adaptive ($\mathrm{F_A}$), Unsustainable ($\mathrm{F_U}$), and Contrasting ($\mathrm{F_C}$).}
\label{fig:supp_Bias_Rec_OM_2}
\end{figure}

\begin{figure}[H]
\centering
\includegraphics[width=\textwidth]{Bias_SSB_OM_2.PNG}
\caption{Relative bias in $SSB$ for region 1, region 2, and global $SSB$ from the first assessment model during the feedback period under the low movement scenario ($\text{OM}_{\text{low}}$). Each point represents the median across the 100 realizations, and error bars represent the interquartile range (25th–75th percentiles). EMs with a grey background indicate that NAA random effects were included, while EMs with a white background indicate that NAA random effects were not included. Three fishing scenarios are included: Adaptive ($\mathrm{F_A}$), Unsustainable ($\mathrm{F_U}$), and Contrasting ($\mathrm{F_C}$).}
\label{fig:supp_Bias_SSB_OM_2}
\end{figure}

\begin{figure}[H]
\centering
\includegraphics[width=\textwidth]{Parameter_Estimates_OM_2.PNG}
\caption{The mean recruitment parameter, recruitment standard deviation and NAA standard deviation from the last assessment model during the feedback period under the low movement scenario ($\text{OM}_{\text{low}}$). Each point represents the median across the 100 realizations, and error bars represent the interquartile range (25th–75th percentiles). EMs with a grey background indicate that NAA random effects were included, while EMs with a white background indicate that NAA random effects were not included. Three fishing scenarios are included: Adaptive ($\mathrm{F_A}$), Unsustainable ($\mathrm{F_U}$), and Contrasting ($\mathrm{F_C}$).}
\label{fig:supp_Parameter_Estimates_OM_2}
\end{figure}

\begin{figure}[H]
\centering
\includegraphics[width=\textwidth]{Catch_G_OM2.PNG}
\caption{Relative difference in annual total catch between the baseline EM ($\text{SpE}_{\scriptscriptstyle \text{NAA,Fix}}$) and other EMs for each fishing history under the low movement scenario ($\text{OM}_{\text{low}}$). Each line represents the median trajectory across the 100 realizations, and the shaded area represents the interquartile range (40th–60th percentiles). The performance of EMs with and without NAA random effects is shown separately for better comparison. Three fishing scenarios are included: Adaptive ($\mathrm{F_A}$), Unsustainable ($\mathrm{F_U}$), and Contrasting ($\mathrm{F_C}$).}
\label{fig:supp_annual_catch_G_OM2}
\end{figure}

\begin{figure}[H]
\centering
\includegraphics[width=\textwidth]{Catch_1_OM2.PNG}
\caption{Relative difference in annual catch in region 1 between the baseline EM ($\text{SpE}_{\scriptscriptstyle \text{NAA,Fix}}$) and other EMs for each fishing history under the low movement scenario ($\text{OM}_{\text{low}}$). Each line represents the median trajectory across the 100 realizations, and the shaded area represents the interquartile range (40th–60th percentiles). The performance of EMs with and without NAA random effects is shown separately for better comparison. Three fishing scenarios are included: Adaptive ($\mathrm{F_A}$), Unsustainable ($\mathrm{F_U}$), and Contrasting ($\mathrm{F_C}$).}
\label{fig:supp_annual_catch_1_OM2}
\end{figure}

\begin{figure}[H]
\centering
\includegraphics[width=\textwidth]{Catch_2_OM2.PNG}
\caption{Relative difference in annual catch in region 2 between the baseline EM ($\text{SpE}_{\scriptscriptstyle \text{NAA,Fix}}$) and other EMs for each fishing history under the low movement scenario ($\text{OM}_{\text{low}}$). Each line represents the median trajectory across the 100 realizations, and the shaded area represents the interquartile range (40th–60th percentiles). The performance of EMs with and without NAA random effects is shown separately for better comparison. Three fishing scenarios are included: Adaptive ($\mathrm{F_A}$), Unsustainable ($\mathrm{F_U}$), and Contrasting ($\mathrm{F_C}$).}
\label{fig:supp_annual_catch_2_OM2}
\end{figure}

\begin{figure}[H]
\centering
\includegraphics[width=\textwidth]{SSB_G_OM2.PNG}
\caption{Relative difference in annual total $SSB$ between the baseline EM ($\text{SpE}_{\scriptscriptstyle \text{NAA,Fix}}$) and other EMs for each fishing history under the low movement scenario ($\text{OM}_{\text{low}}$). Each line represents the median trajectory across the 100 realizations, and the shaded area represents the interquartile range (40th–60th percentiles). The performance of EMs with and without NAA random effects is shown separately for better comparison. Three fishing scenarios are included: Adaptive ($\mathrm{F_A}$), Unsustainable ($\mathrm{F_U}$), and Contrasting ($\mathrm{F_C}$).}
\label{fig:supp_SSB_G_OM2}
\end{figure}

\begin{figure}[H]
\centering
\includegraphics[width=\textwidth]{SSB_1_OM2.PNG}
\caption{Relative difference in annual $SSB$ for region 1 between the baseline EM ($\text{SpE}_{\scriptscriptstyle \text{NAA,Fix}}$) and other EMs for each fishing history under the low movement scenario ($\text{OM}_{\text{low}}$). Each line represents the median trajectory across the 100 realizations, and the shaded area represents the interquartile range (40th–60th percentiles). The performance of EMs with and without NAA random effects is shown separately for better comparison. Three fishing scenarios are included: Adaptive ($\mathrm{F_A}$), Unsustainable ($\mathrm{F_U}$), and Contrasting ($\mathrm{F_C}$).}
\label{fig:supp_SSB_1_OM2}
\end{figure}

\begin{figure}[H]
\centering
\includegraphics[width=\textwidth]{SSB_2_OM2.PNG}
\caption{Relative difference in annual $SSB$ for region 2 between the baseline EM ($\text{SpE}_{\scriptscriptstyle \text{NAA,Fix}}$) and other EMs for each fishing history under the low movement scenario ($\text{OM}_{\text{low}}$). Each line represents the median trajectory across the 100 realizations, and the shaded area represents the interquartile range (40th–60th percentiles). The performance of EMs with and without NAA random effects is shown separately for better comparison. Three fishing scenarios are included: Adaptive ($\mathrm{F_A}$), Unsustainable ($\mathrm{F_U}$), and Contrasting ($\mathrm{F_C}$).}
\label{fig:supp_SSB_2_OM2}
\end{figure}

\begin{figure}[H]
\centering
\includegraphics[width=\textwidth]{Performance_first_5_years_relative_OM_2.PNG}
\caption{Relative difference in catch, $F$, and $SSB$ between the baseline EM ($\text{SpE}_{\scriptscriptstyle \text{NAA,Fix}}$) and other EMs over the first 5 years of the feedback period under the low movement scenario ($\text{OM}_{\text{low}}$). Each point represents the median across the 100 realizations, and error bars represent the interquartile range (25th–75th percentiles). EMs with a grey background indicate that NAA random effects were included, while EMs with a white background indicate that NAA random effects were not included. Three fishing scenarios are included: Adaptive ($\mathrm{F_A}$), Unsustainable ($\mathrm{F_U}$), and Contrasting ($\mathrm{F_C}$).}
\label{fig:supp_first_5yr_results_OM2}
\end{figure}

\begin{figure}[H]
\centering
\includegraphics[width=\textwidth]{Performance_last_5_years_relative_OM_2.PNG}
\caption{Relative difference in catch, $F$, and $SSB$ between the baseline EM ($\text{SpE}_{\scriptscriptstyle \text{NAA,Fix}}$) and other EMs over the last 5 years of the feedback period under the low movement scenario ($\text{OM}_{\text{low}}$). Each point represents the median across the 100 realizations, and error bars represent the interquartile range (25th–75th percentiles). EMs with a grey background indicate that NAA random effects were included, while EMs with a white background indicate that NAA random effects were not included. Three fishing scenarios are included: Adaptive ($\mathrm{F_A}$), Unsustainable ($\mathrm{F_U}$), and Contrasting ($\mathrm{F_C}$).}
\label{fig:supp_last_5yr_results_OM2}
\end{figure}

\begin{figure}[H]
\centering
\includegraphics[width=\textwidth]{Catch_variation_OM2.PNG}
\caption{Average annual catch variation (AACV) of each EM at both regional and global scales under the low movement scenario ($\text{OM}_{\text{low}}$). Each point represents the median across the 100 realizations (one value per realization), and error bars represent the interquartile range (25th–75th percentiles). EMs with a grey background indicate that NAA random effects were included, while EMs with a white background indicate that NAA random effects were not included. Three fishing scenarios are included: Adaptive ($\mathrm{F_A}$), Unsustainable ($\mathrm{F_U}$), and Contrasting ($\mathrm{F_C}$).}
\label{fig:supp_aacv_OM2}
\end{figure}

\begin{figure}[H]
\centering
\includegraphics[width=\textwidth]{Probability_Overfished_Overfishing_OM2_MSY.PNG}
\caption{Probability of $SSB < SSB_{MSY}$ and $F > F_{MSY}$ for each EM over the 30-year feedback period under the low movement scenario ($\text{OM}_{\text{low}}$). Each point represents the median across the 100 realizations (one value per realization), and error bars represent the interquartile range (25th–75th percentiles). EMs with a grey background indicate that NAA random effects were included, while EMs with a white background indicate that NAA random effects were not included. Three fishing scenarios are included: Adaptive ($\mathrm{F_A}$), Unsustainable ($\mathrm{F_U}$), and Contrasting ($\mathrm{F_C}$).}
\label{fig:supp_Overfished_OM2_MSY}
\end{figure}

\begin{figure}[H]
\centering
\includegraphics[width=\textwidth]{Probability_Overfished_Overfishing_OM2_SPR.PNG}
\caption{Probability of $SSB < SSB_{40\%}$ and $F > F_{40\%}$ for each EM over the 30-year feedback period under the low movement scenario ($\text{OM}_{\text{low}}$). Each point represents the median across the 100 realizations (one value per realization), and error bars represent the interquartile range (25th–75th percentiles). EMs with a grey background indicate that NAA random effects were included, while EMs with a white background indicate that NAA random effects were not included. Three fishing scenarios are included: Adaptive ($\mathrm{F_A}$), Unsustainable ($\mathrm{F_U}$), and Contrasting ($\mathrm{F_C}$).}
\label{fig:supp_Overfished_OM2_SPR}
\end{figure}

\begin{figure}[H]
\centering
\includegraphics[width=\textwidth]{Population_Status_OM_2_MSY.PNG}
\caption{Terminal-year status of $F$ and $SSB$ at the global scale for each EM under the low movement scenario ($\text{OM}_{\text{low}}$). Each point represents the median across the 100 realizations, and error bars represent the interquartile range (25th–75th percentiles). The vertical dashed line indicates the threshold of overfished ($SSB_{T}/SSB_{MSY_{T}} < 0.5$), and the horizontal dashed line indicates the threshold of overfishing ($F_{T}/F_{MSY_{T}} > 1$). The performance of EMs with and without NAA random effects is shown separately for better comparison. Three fishing scenarios are included: Adaptive ($\mathrm{F_A}$), Unsustainable ($\mathrm{F_U}$), and Contrasting ($\mathrm{F_C}$).}
\label{fig:supp_population_status2_MSY}
\end{figure}

\begin{figure}[H]
\centering
\includegraphics[width=\textwidth]{Population_Status_OM_2_SPR.PNG}
\caption{Terminal-year status of $F$ and $SSB$ at the global scale for each EM under the low movement scenario ($\text{OM}_{\text{low}}$). Each point represents the median across the 100 realizations, and error bars represent the interquartile range (25th–75th percentiles). The vertical dashed line indicates the threshold of overfished ($SSB_{T}/SSB_{40\%_{T}} < 0.5$), and the horizontal dashed line indicates the threshold of overfishing ($F_{T}/F_{40\%_{T}} > 1$). The performance of EMs with and without NAA random effects is shown separately for better comparison. Three fishing scenarios are included: Adaptive ($\mathrm{F_A}$), Unsustainable ($\mathrm{F_U}$), and Contrasting ($\mathrm{F_C}$).}
\label{fig:supp_population_status2_SPR}
\end{figure}

\begin{figure}[H]
\centering
\includegraphics[width=\textwidth,height=\textheight]{Holistic_OM_2.PNG}
\caption{Overall performance of each EM at the global scale under the low movement scenario ($\text{OM}_{\text{low}}$), including the following relative performance metrics: 1) long-term catch ($\text{Catch}_{\text{long}}$); 2) short-term catch ($\text{Catch}_{\text{short}}$); 3) long-term $SSB$ ($SSB_{\text{long}}$); 4) short-term $SSB$ ($SSB_{\text{short}}$); 5) probability of $F < F_{MSY}$ ($\text{P}_{F<F_{MSY}}$); 6) probability of $SSB > SSB_{MSY}$ ($\text{P}_{SSB>SSB_{MSY}}$); 7) overfishing status in the terminal year ($\text{OFG}_{\text{terminal}}$); 8) overfished status in the terminal year ($\text{OFD}_{\text{terminal}}$); and 9) average annual catch variation (AACV). Bars represent the median score across the 100 realizations, and error bars represent the interquartile range (25th–75th percentiles). The number displayed in each panel represents the total score (sum of all medians for that EM) for the corresponding EM under that fishing scenario, with higher totals indicating better overall performance. The black dashed line separates the performance of each EM. EMs with a grey background indicate that NAA random effects were included, while EMs with a white background indicate that NAA random effects were not included. Three fishing scenarios are included: Adaptive ($\mathrm{F_A}$), Unsustainable ($\mathrm{F_U}$), and Contrasting ($\mathrm{F_C}$).}
\label{fig:supp_holistic2}
\end{figure}

\begin{figure}[H]
\centering
\includegraphics[width=\textwidth,height=\textheight]{Holistic_OM_R1_2.PNG}
\caption{Overall performance of each EM in region 1 under the low movement scenario ($\text{OM}_{\text{low}}$), including the following relative performance metrics: 1) long-term catch ($\text{Catch}_{\text{long}}$); 2) short-term catch ($\text{Catch}_{\text{short}}$); 3) long-term $SSB$ ($SSB_{\text{long}}$); 4) short-term $SSB$ ($SSB_{\text{short}}$); 5) probability of $F < F_{MSY}$ ($\text{P}_{F<F_{MSY}}$); 6) probability of $SSB > SSB_{MSY}$ ($\text{P}_{SSB>SSB_{MSY}}$); 7) overfishing status in the terminal year ($\text{OFG}_{\text{terminal}}$); 8) overfished status in the terminal year ($\text{OFD}_{\text{terminal}}$); and 9) average annual catch variation (AACV). Bars represent the median score across the 100 realizations, and error bars represent the interquartile range (25th–75th percentiles). The number displayed in each panel represents the total score (sum of all medians for that EM) for the corresponding EM under that fishing scenario, with higher totals indicating better overall performance. The black dashed line separates the performance of each EM. EMs with a grey background indicate that NAA random effects were included, while EMs with a white background indicate that NAA random effects were not included. Three fishing scenarios are included: Adaptive ($\mathrm{F_A}$), Unsustainable ($\mathrm{F_U}$), and Contrasting ($\mathrm{F_C}$).}
\label{fig:supp_holistic2_R1}
\end{figure}

\begin{figure}[H]
\centering
\includegraphics[width=\textwidth,height=\textheight]{Holistic_OM_R2_2.PNG}
\caption{Overall performance of each EM in region 2 under the low movement scenario ($\text{OM}_{\text{low}}$), including the following relative performance metrics: 1) long-term catch ($\text{Catch}_{\text{long}}$); 2) short-term catch ($\text{Catch}_{\text{short}}$); 3) long-term $SSB$ ($SSB_{\text{long}}$); 4) short-term $SSB$ ($SSB_{\text{short}}$); 5) probability of $F < F_{MSY}$ ($\text{P}_{F<F_{MSY}}$); 6) probability of $SSB > SSB_{MSY}$ ($\text{P}_{SSB>SSB_{MSY}}$); 7) overfishing status in the terminal year ($\text{OFG}_{\text{terminal}}$); 8) overfished status in the terminal year ($\text{OFD}_{\text{terminal}}$); and 9) average annual catch variation (AACV). Bars represent the median score across the 100 realizations, and error bars represent the interquartile range (25th–75th percentiles). The number displayed in each panel represents the total score (sum of all medians for that EM) for the corresponding EM under that fishing scenario, with higher totals indicating better overall performance. The black dashed line separates the performance of each EM. EMs with a grey background indicate that NAA random effects were included, while EMs with a white background indicate that NAA random effects were not included. Three fishing scenarios are included: Adaptive ($\mathrm{F_A}$), Unsustainable ($\mathrm{F_U}$), and Contrasting ($\mathrm{F_C}$).}
\label{fig:supp_holistic2_R2}
\end{figure}

\renewcommand{\thefigure}{\arabic{figure}}

\end{document}
